%! TEX program = xelatex
%! TEX encoding = UTF-8 Unicode

\documentclass[internationalmaster]{hquThesis}
\usepackage{lineno}
\usepackage{amsmath,amsfonts,amssymb,textcomp,ulem}
\usepackage{booktabs,tabularray,float}
\usepackage{listings}
\graphicspath{{figures/}}



\hypersetup{colorlinks=true,linkcolor=black,citecolor=black}

\begin{document}
\makecover
% NOTE: 正式文档请取消下面两行注释
% \include{decision}
% \include{declaration}

\frontmatter{
    \include{abstract}
    \tableofcontents
}

\mainmatter
\chapter{前期准备}
\section{封面关键字设置}
\section{Git相关}
\begin{enumerate}
    \item 通过fork建仓，具体步骤可参考\href{https://docs.github.com/en/repositories/creating-and-managing-repositories/creating-a-template-repository}{这里}。
    \item 设置 .gitignore 文件(已配置)。
\item 设置临时文件(.aux、.bbl等)和编译文件(.pdf)位置，具体可参考\href{https://github.com/James-Yu/LaTeX-Workshop/wiki/Compile#latex-tools}{LaTeX-Workshop帮助文档}(设置\%TMPDIR\%、\%OUTDIR\%)。
\end{enumerate}

\section{VS code配置}
\subsection*{配置Latex}
\href{https://github.com/James-Yu/LaTeX-Workshop/wiki/Compile#latex-tools}{LaTeX-Workshop配置}
\subsection*{LaTeX workshop: }

正向搜索(forward research)从代码处定位至PDF文档处：Ctrl+alt+j

反向搜索(inverse research)从PDF文档准确定位至LaTeX代码处：Ctrl+鼠标左键

正反向搜索具体教程可参考\href{https://github.com/James-Yu/LaTeX-Workshop/wiki/View#synctex}{这里}。
\subsection*{Git Graph}
点击左侧源代码管理插件-点击左上角View Git Graph(git log)-查看修改记录

\subsection*{Todo Tree}
帮助管理项目中的TODO注释和其他标记，\href{https://blog.csdn.net/m0_37586991/article/details/104239568}{Todo Tree安装步骤}。

\subsection*{Zotero LaTeX}
Zotero管理参考文献
\section{文件、label 命名规则}
Latex路径中文件只允许以英文命名

\subsection*{公式}
\begin{lstlisting}[language=TeX]
    \label{CH_name}{formula_name}
\end{lstlisting}

\subsection*{图}
\begin{lstlisting}
    \label{CH_name}{fig_name}
\end{lstlisting}

\subsection*{表}
\begin{lstlisting}
    \label{CH_name}{table_name}
\end{lstlisting}



\chapter{剪切自锁问题}

\section{剪切自锁现象的物理机制}

核心要点：从总势能泛函出发，揭示厚跨比变薄时横向剪切能量项退化为"罚函数"导致结构刚度被高估的数学本质。

% TODO: 补充剪切自锁物理机制的详细推导与论述

\section{传统有限元免剪切自锁的方案}

核心要点：概述传统有限元通过数值积分或多变量变分手段，合理放宽虚假剪切约束的三大主流技术演进路线。

\subsection{减缩积分方案}

核心要点：剖析降低剪切刚度矩阵数值积分阶次、释放奇异值以软化数值刚度的积分阶次效应。

% TODO: 补充减缩积分方案的原理与算例

\subsection{应变空间重构方案}

核心要点：说明利用特定几何采样点进行剪切应变重构，从而完全消除寄生剪切力的基本原理。引入独立插值场的多变量变分原理，放宽位移场与剪切约束耦合度的解锁机制。

% TODO: 补充应变空间重构方案（ANS、MITC等）的详细论述

\section{LBB 稳定性条件的最优约束比}

核心要点：归纳上述解锁策略"合理放宽约束"的变分共性，作为构建本文混合离散格式的理论基础。推导混合离散 LBB 不等式，阐明其在完全消除自锁（免锁死）与抑制虚假零能模态（沙漏模态）之间的代数平衡机制。

% TODO: 补充 LBB 条件与约束比的数学推导

\section{小结}

% TODO: 第二章小结

\chapter{自由振动剪切自锁有限元无网格混合离散分析}

% 第三章：自由振动分析公式推导
% 包含：混合离散动力学方程、能量方程、弱形式

本章针对 Mindlin 中厚板的自由振动问题，基于第二章建立的运动学方程，推导混合离散格式下的动力学控制方程与广义特征值问题。首先推导板的动能变分与应变能变分，通过哈密顿原理建立自由振动的虚功方程；其次进行离散化，得到材料刚度矩阵与质量矩阵的各子块表达式；最后讨论混合离散格式的稳定性条件。

\section{混合离散动力学方程与特征值问题}

考虑 Mindlin 中厚板的自由振动问题，求解目标为结构的固有频率 $\omega$ 和固有振型。自由振动对应的广义特征值问题为：
\begin{equation}
(\mathbf{K} - \omega^2 \mathbf{M})\mathbf{D} = \mathbf{0}
\end{equation}
式中，$\mathbf{K}$ 为弹性刚度矩阵，$\mathbf{M}$ 为质量矩阵，$\mathbf{D}$ 为节点自由度向量。

\section{振动分析能量方程}

由第二章建立的 Mindlin 板位移场、应变-位移关系及材料本构可知，本章直接在此基础上进行动力分析。首先推导板的动能变分与应变能变分，再通过哈密顿原理建立自由振动的虚功方程。板的动能变分为：
\begin{equation}
\delta T = \int_V \rho \ddot{u}_i \delta u_i\, dV
\end{equation}
将位移场代入，利用中面位置 $x_3=0$ 的对称性以及密度均匀假设，定义惯性系数
\begin{equation}
I_0 = \int_{-h/2}^{h/2} \rho\, dx_3 = \rho h, \qquad I_2 = \int_{-h/2}^{h/2} \rho x_3^2\, dx_3 = \frac{\rho h^3}{12}
\end{equation}
由于中面对称性，$\int_{-h/2}^{h/2} \rho x_3 dx_3 = 0$，可得动能变分的中面积分形式：
\begin{equation}
\delta T = \int_\Omega \left(I_0 \ddot{w}\,\delta w + I_2 \ddot{\psi}_\alpha \delta\psi_\alpha\right) d\Omega
\end{equation}

板的应变能变分为：
\begin{equation}
\delta U = \int_V \sigma_{ij} \delta\varepsilon_{ij}\, dV
\end{equation}
内虚功密度 $\sigma_{ij}\delta\varepsilon_{ij}$ 包含面内分量与横向分量（$\sigma_{33}=0$ 不计）：
\begin{equation}
\sigma_{ij}\delta\varepsilon_{ij} = \sigma_{\alpha\beta}\delta\varepsilon_{\alpha\beta} + 2\sigma_{\alpha 3}\delta\varepsilon_{\alpha 3}
\end{equation}
利用 Mindlin 位移场对应的应变变分 $\delta\varepsilon_{\alpha\beta} = x_3\delta\kappa_{\alpha\beta}$，$\delta\varepsilon_{\alpha 3} = \frac{1}{2}\delta\gamma_\alpha$，代入并对厚度积分，定义弯矩合力与剪力合力
\begin{equation}
M_{\alpha\beta} \equiv \int_{-h/2}^{h/2} x_3 \sigma_{\alpha\beta}\, dx_3 = \frac{h^3}{12} C_{\alpha\beta\gamma\delta}\kappa_{\gamma\delta}, \qquad Q_\alpha \equiv \int_{-h/2}^{h/2} \sigma_{\alpha 3}\, dx_3 = \kappa G h \gamma_\alpha
\end{equation}
由此，内虚功可写为
\begin{equation}
\delta U = \int_\Omega \left(M_{\alpha\beta}\delta\kappa_{\alpha\beta} + Q_\alpha \delta\gamma_\alpha\right) d\Omega
\end{equation}

根据哈密顿原理，系统在任意时间区间 $[t_1, t_2]$ 内的动能变分与应变能变分满足：
\begin{equation}
\int_{t_1}^{t_2} (\delta T - \delta U)\, dt = 0,
\end{equation}
对任意虚位移 $\delta u_i$，可得弹性动力学的虚功方程：
\begin{equation}
\int_V (\sigma_{ij}\delta\varepsilon_{ij} + \rho \ddot{u}_i \delta u_i)\, dV = 0.
\end{equation}
板的体积积分可分解为中面积分与厚度积分：$\int_V dV = \int_\Omega \int_{-h/2}^{h/2} dx_3\, d\Omega$，其中 $\Omega$ 为板的中面。将内虚功与惯性项代入，得到 Mindlin 板自由振动的虚功方程：
\begin{equation}
\int_\Omega \left[M_{\alpha\beta}\delta\kappa_{\alpha\beta} + Q_\alpha \delta\gamma_\alpha - I_0 \ddot{w}\,\delta w - I_2 \ddot{\psi}_\alpha \delta \psi_\alpha\right] d\Omega = 0.
\end{equation}

\section{弱形式}

对上述虚功方程各变量取变分并令其为零，将离散场代入后得到广义特征值方程。位移场和转角场的形函数插值为：
\begin{equation}
w(\mathbf{x}) = N_I w_I, \qquad \varphi_\alpha(\mathbf{x}) = N_I \varphi_{\alpha I}
\end{equation}
对空间坐标求偏导，微分算子直接作用在形函数上：$w_{,\alpha} = N_{I,\alpha}w_I$，$\varphi_{\alpha,\beta} = N_{I,\beta}\varphi_\alpha$。虚位移做同阶离散近似。

材料刚度矩阵各子块为
\begin{equation}
K_{IJ}^{ww} = \int_\Omega \kappa G h\, N_{I,\alpha} N_{J,\alpha}\, d\Omega, \qquad K_{I,J\gamma}^{w\varphi} = -\int_\Omega \kappa G h\, N_{I,\gamma} N_J\, d\Omega
\end{equation}
\begin{equation}
K_{I\alpha,J\gamma}^{\varphi\varphi} = \int_\Omega \left(N_{I,\beta} D_{\alpha\beta\gamma\eta} N_{J,\eta} + \kappa G h\, N_I N_J \delta_{\alpha\gamma}\right) d\Omega
\end{equation}

质量矩阵展开后的对称矩阵形式为
\begin{equation}
\mathbf{M} = \begin{bmatrix}
M^{ww} & \mathbf{0} & \mathbf{0} \\
\mathbf{0} & M^{\varphi_1\varphi_1} & \mathbf{0} \\
\mathbf{0} & \mathbf{0} & M^{\varphi_2\varphi_2}
\end{bmatrix}
\end{equation}
其中横向挠度对应的质量子块为：
\begin{equation}
M_{IJ}^{ww} = \int_\Omega I_0 N_I N_J\, d\Omega = \int_\Omega \rho h\, N_I N_J\, d\Omega
\end{equation}
转角自由度对应的质量子块为：
\begin{equation}
M_{I\alpha,J\beta}^{\varphi\varphi} = \int_\Omega I_2 \delta_{\alpha\beta} N_I N_J\, d\Omega = \int_\Omega \frac{\rho h^3}{12} \delta_{\alpha\beta} N_I N_J\, d\Omega
\end{equation}

由于线性问题中面内与弯曲不耦合，全局特征值方程呈现块对角结构：
\begin{equation}
\begin{bmatrix}
\mathbf{K}_{ww} & \mathbf{K}_{w\varphi} \\
\mathbf{K}_{\varphi w} & \mathbf{K}_{\varphi\varphi}
\end{bmatrix}
\begin{bmatrix}
\mathbf{W} \\
\boldsymbol{\Psi}
\end{bmatrix}
= \omega^2
\begin{bmatrix}
\mathbf{M}_{ww} & \mathbf{0} \\
\mathbf{0} & \mathbf{M}_{\varphi\varphi}
\end{bmatrix}
\begin{bmatrix}
\mathbf{W} \\
\boldsymbol{\Psi}
\end{bmatrix}
\end{equation}
此即 Mindlin 板自由振动的全局离散特征值方程。求解广义特征值问题 $(\mathbf{K} - \omega^2\mathbf{M})\mathbf{D} = \mathbf{0}$，可得到自然频率 $\omega$ 和对应模态。

格式稳定性方面，振动问题中，质量矩阵 $\mathbf{M}$ 为对称正定。为保证自由振动特征值问题的适定性，混合离散格式需满足 LBB 稳定性条件。本文采用的剪切应力场插值阶数经约束比分析验证，可保证 inf-sup 常数有正下界，从而在薄板极限下避免剪切自锁并维持刚度矩阵的正定性。当板厚趋近于零时，最低阶频率收敛于 Kirchhoff 薄板理论的解析解。

% 方板自由振动数值算例
% 包含：四边固支(CCCC)、四边简支(SSSS)

\section{数值算例}

通过数值算例系统检验不同厚跨比与边界条件下，混合格式在自由振动分析中的计算精度。

自由振动算例采用与屈曲算例一致的无量纲材料参数，弹性模量 $E=1.0$，泊松比 $\nu=0.3$，密度 $\rho=1.0$，板厚 $h=10^{-2}$，板边长 $a=b=1.0$，剪切修正系数 $\kappa=5/6$。网格采用结构化四边形单元（Quad4），规模为 $25\times25$ 至 $35\times35$。

\subsection{方板边界条件自由振动分析}

% TODO: 补充四边固支方板自由振动算例数据与分析

针对四边简支边界进行计算。

% TODO: 补充四边简支方板自由振动算例数据与分析


\section{小结}

% TODO: 第三章小结

\chapter{屈曲分析剪切自锁有限元无网格混合离散分析}

% 第四章：屈曲分析公式推导
% 包含：混合离散屈曲控制方程、能量方程、弱形式

本章针对 Mindlin 中厚板的屈曲问题，基于第二章建立的运动学方程，推导混合离散格式下的屈曲控制方程与广义特征值问题。首先采用 Green-Lagrange 应变张量推导屈曲问题的虚功方程；其次进行离散化，得到材料刚度矩阵与几何刚度矩阵的各子块表达式；最后讨论混合离散格式的稳定性条件。

\section{混合离散屈曲控制方程与特征值问题}

考虑 Mindlin 中厚板的屈曲问题。板在面内压力作用下，当荷载增大至临界值时，原始平面内平衡状态失去稳定性，板突然发生横向挠度（屈曲分岔）。与强度失效不同，屈曲失效通常发生在应力远低于材料屈服强度时，具有突发性。求解目标为临界荷载系数 $\lambda$ 和屈曲模态，对应的广义特征值问题为：
\begin{equation}
(\mathbf{K} + \lambda \mathbf{K}_G)\boldsymbol{\phi} = \mathbf{0}
\end{equation}
式中，$\mathbf{K}$ 为弹性刚度矩阵，$\mathbf{K}_G$ 为几何刚度矩阵，$\lambda$ 为屈曲荷载因子，$\boldsymbol{\phi}$ 为屈曲模态。临界屈曲荷载为 $P_{cr} = \lambda_{cr} P_{ref}$，其中 $P_{ref}$ 为参考预加荷载，$\lambda_{cr}$ 为最小正特征值。

面内位移变分 $u_\alpha$ 不会与预应力度 $\sigma_{\alpha\beta}$ 产生导致屈曲的几何力矩，因此几何刚度矩阵中与 $u_\alpha$ 对应的行列皆为零。实际控制板分岔屈曲的广义特征值问题子块矩阵方程为：
\begin{equation}
(\mathbf{k}_b + \lambda \mathbf{k}_g)\mathbf{b} = \mathbf{0}
\end{equation}

\section{屈曲分析能量方程}

由第二章建立的 Mindlin 板位移场可知，屈曲分析需考虑几何非线性，采用完整的 Green-Lagrange 应变张量：
\begin{equation}
\varepsilon_{ij} = \frac{1}{2}(\bar{u}_{i,j} + \bar{u}_{j,i} + \bar{u}_{k,i}\bar{u}_{k,j})
\end{equation}

将位移场代入，面内应变分量 $\varepsilon_{\alpha\beta}$ 按 $x_3$ 的幂次分项归类为：
\begin{equation}
\begin{aligned}
\varepsilon_{\alpha\beta} =& \underbrace{\frac{1}{2}(u_{\alpha,\beta} + u_{\beta,\alpha} + u_{\gamma,\alpha}u_{\gamma,\beta})}_{\text{薄膜应变}} + \underbrace{\frac{1}{2}w_{,\alpha}w_{,\beta}}_{\text{von K\'arm\'an 项}} \\
&- x_3\underbrace{\frac{1}{2}(\varphi_{\alpha,\beta} + \varphi_{\beta,\alpha} + u_{\gamma,\alpha}\varphi_{\gamma,\beta} + \varphi_{\gamma,\alpha}u_{\gamma,\beta})}_{\text{广义曲率}} \\
&+ x_3^2\underbrace{\frac{1}{2}\varphi_{\gamma,\alpha}\varphi_{\gamma,\beta}}_{\text{高阶旋转项}}
\end{aligned}
\end{equation}

横向剪切应变 $\varepsilon_{\alpha 3}$ 和厚度方向法向应变 $\varepsilon_{33}$ 分别为：
\begin{equation}
\varepsilon_{\alpha 3} = \frac{1}{2}(w_{,\alpha} - \varphi_\alpha - u_{\beta,\alpha}\varphi_\beta) + \frac{x_3}{2}\varphi_{\beta,\alpha}\varphi_\beta, \qquad \varepsilon_{33} = \frac{1}{2}\varphi_\alpha\varphi_\alpha
\end{equation}
上述各应变分量的完整逐项展开过程参见附录 B。

根据虚功原理，板的内虚功等于外虚功：
\begin{equation}
\int_{-h/2}^{h/2} \varepsilon_{ij}\sigma_{ij}\, dx_3\, d\Omega = \delta W_{ext}
\end{equation}
将应变张量按面内分量、横向剪切分量和厚度方向分量拆分为三部分积分：
\begin{equation}
\int_\Omega\int_{-h/2}^{h/2} \varepsilon_{\alpha\beta}\sigma_{\alpha\beta}\, dx_3 d\Omega + 2\int_\Omega\int_{-h/2}^{h/2} \varepsilon_{\alpha 3}\sigma_{\alpha 3}\, dx_3 d\Omega + \int_\Omega\int_{-h/2}^{h/2} \varepsilon_{33}\sigma_{33}\, dx_3 d\Omega = \delta W_{ext}
\end{equation}

第一项为面内应变项。将总应力 $\sigma_{\alpha\beta}$ 拆分为预应力度与屈曲增量应力：
\begin{equation}
\sigma_{\alpha\beta} = \bar{\sigma}_{\alpha\beta} + \sigma_{\alpha\beta}^{inc}
\end{equation}
其中，预应力度 $\bar{\sigma}_{\alpha\beta}$ 沿厚度均匀分布（与 $x_3$ 无关）；增量应力 $\sigma_{\alpha\beta}^{inc}$ 随厚度弹性变化。将应变展开式代入并完全展开，面内应变分量与应力的乘积共含 8 项：(1) 公式(4.4)中薄膜应变部分与预应力度 $\bar{\sigma}_{\alpha\beta}$ 的乘积；(2) 薄膜应变部分与增量应力 $\sigma_{\alpha\beta}^{inc}$ 的乘积；(3) 广义曲率项 $-x_3\kappa_{\alpha\beta}$ 与预应力度的乘积；(4) 广义曲率项与增量应力的乘积；(5) von K\'arm\'an 项 $\frac{1}{2}w_{,\alpha}w_{,\beta}$ 与预应力度的乘积；(6) von K\'arm\'an 项与增量应力的乘积；(7) 高阶旋转项 $x_3^2\eta_{\alpha\beta}$ 与预应力度的乘积；(8) 高阶旋转项与增量应力的乘积。依据小变形假设和屈曲分岔点的特性，第 1、2 项涉及膜应变，其一阶变分为零；第 6、8 项为增量应力与二阶位移梯度的乘积，属三阶小量舍去；第 3 项因预应力度沿厚度均匀分布与 $x_3$ 无关，$\int x_3 dx_3 = 0$ 恒消去。保留第 4、5、7 项，对厚度积分后定义弯矩合力：
\begin{equation}
M_{\alpha\beta} = \int_{-h/2}^{h/2} x_3 \sigma_{\alpha\beta}^{inc}\, dx_3
\end{equation}
第一项积分的最终结果为：
\begin{equation}
\int_\Omega\int_{-h/2}^{h/2} \varepsilon_{\alpha\beta}\sigma_{\alpha\beta}\, dx_3 d\Omega = \int_\Omega \kappa_{\alpha\beta}M_{\alpha\beta}\, d\Omega + \int_\Omega w_{,\alpha}(h\bar{\sigma}_{\alpha\beta})w_{,\beta}\, d\Omega + \int_\Omega \varphi_{\gamma,\alpha}\left(\frac{h^3}{12}\bar{\sigma}_{\alpha\beta}\right)\varphi_{\gamma,\beta}\, d\Omega
\end{equation}

第二项为横向剪切项。对于纯面内加载的板，预屈曲横向剪应力 $\bar{\sigma}_{\alpha 3} = 0$，故 $\sigma_{\alpha 3} = \sigma_{\alpha 3}^{inc}$。由虎克定律，引入剪切修正系数 $\kappa$ 后：
\begin{equation}
\sigma_{\alpha 3}^{inc} = \kappa G(w_{,\alpha} - \varphi_\alpha)
\end{equation}
代入并对厚度积分（$\int_{-h/2}^{h/2} dx_3 = h$），注意工程剪切应变 $\gamma_\alpha = 2\varepsilon_{\alpha 3}$，第二项的积分结果为：
\begin{equation}
2\int_\Omega\int_{-h/2}^{h/2} \varepsilon_{\alpha 3}\sigma_{\alpha 3}\, dx_3 d\Omega = \int_\Omega \kappa G h\,(w_{,\alpha} - \varphi_\alpha)(w_{,\alpha} - \varphi_\alpha)\, d\Omega
\end{equation}

第三项为厚度方向法向应变项，在平面应力假设 $\sigma_{33}=0$ 下贡献为零。

将三项合并，得到屈曲问题的完整虚功方程：
\begin{equation}
\begin{aligned}
\delta W_{ext} =& \int_\Omega \kappa_{\alpha\beta}M_{\alpha\beta}\, d\Omega + \int_\Omega \kappa G h\,(w_{,\alpha} - \varphi_\alpha)(w_{,\alpha} - \varphi_\alpha)\, d\Omega \\
&+ \int_\Omega h\bar{\sigma}_{\alpha\beta} w_{,\alpha}w_{,\beta}\, d\Omega + \int_\Omega \frac{h^3}{12}\bar{\sigma}_{\alpha\beta} \varphi_{\gamma,\alpha}\varphi_{\gamma,\beta}\, d\Omega
\end{aligned}
\end{equation}
其中前两项为弹性内力虚功（弯曲与剪切），后两项为预应力度对应的几何刚度虚功。当几何刚度项的绝对值等于弹性刚度项时，系统总刚度矩阵奇异，达到临界屈曲状态。

\section{弱形式}

对上述屈曲虚功方程各变量取变分并令其为零，将离散场代入后得到离散形式的广义特征值方程。位移场和转角场的形函数插值为：
\begin{equation}
w(\mathbf{x}) = N_I w_I, \qquad \varphi_\alpha(\mathbf{x}) = N_I \varphi_{\alpha I}
\end{equation}
对空间坐标求偏导，微分算子直接作用在形函数上：$w_{,\alpha} = N_{I,\alpha}w_I$，$\varphi_{\alpha,\beta} = N_{I,\beta}\varphi_\alpha$。虚位移做同阶离散近似。

材料刚度矩阵子块展开如下：面外材料刚度矩阵 $k_b$ 由弯曲刚度（Bending）与横向剪切刚度（Shear）共同组成。对应自由度 $[w,\varphi_1,\varphi_2]^T$，展开后的对称矩阵形式如下：
\begin{equation}
k_b = \begin{bmatrix}
k_{ww} & k_{w\varphi_1} & k_{w\varphi_2} \\
k_{\varphi_1 w} & k_{\varphi_1\varphi_1} & k_{\varphi_1\varphi_2} \\
k_{\varphi_2 w} & k_{\varphi_2\varphi_1} & k_{\varphi_2\varphi_2}
\end{bmatrix}
\end{equation}

各子块的积分表达式为：对于 $k_{ww}$，
\begin{equation}
k_{ww} = \int \kappa Gh(w_{,\alpha}w_{,\alpha})d\Omega
\end{equation}
\begin{equation}
k_{ww} = \int \kappa Gh(N_{I,\alpha}w_I\cdot N_{J,\alpha}w_J)d\Omega
\end{equation}
\begin{equation}
k_{ww} = w_I\cdot\left[\int \kappa Gh(N_{I,\alpha}N_{J,\alpha})d\Omega\right]\cdot w_J
\end{equation}
\begin{equation}
\boxed{K_{IJ}^{ww} = \int \kappa Gh(N_{I,\alpha}N_{J,\alpha})d\Omega}
\end{equation}

对于 $k_{w\varphi_\gamma}$，
\begin{equation}
k_{w\varphi_\gamma} = -\int \kappa Gh(w_{,\alpha}\varphi_\alpha)d\Omega
\end{equation}
\begin{equation}
k_{IJ}^{w\varphi_\gamma} = -\int \kappa Gh(N_{I,\alpha}w_I\cdot N_{\alpha\gamma}N_J\varphi_\gamma)d\Omega
\end{equation}
\begin{equation}
k_{IJ}^{w\varphi_\gamma} = w_I\cdot\left[-\int \kappa Gh(N_{I,\gamma}N_J)d\Omega\right]\cdot\varphi_\gamma
\end{equation}
\begin{equation}
\boxed{k_{IJ}^{w\varphi_\gamma} = -\int \kappa Gh(N_{I,\gamma}N_J)d\Omega}
\end{equation}

对于 $k_{\varphi_\alpha\varphi_\gamma}$，
弯矩合力与曲率的关系为 $M_{\alpha\beta}=D_{\alpha\beta\gamma}\kappa_\gamma$，$\kappa_{\alpha\beta}=\frac{1}{2}(\varphi_{\alpha,\beta}+\varphi_{\beta,\alpha})$。
\begin{equation}
k_{\alpha\gamma}^{\varphi\varphi} = \int(\phi_{\alpha,\beta}D_{\alpha\beta\gamma\eta}\phi_{\gamma,\eta} + \kappa Gh\phi_\alpha\phi_\alpha)d\Omega
\end{equation}
\begin{equation}
k_{I\alpha,J\gamma}^{\varphi\varphi} = \phi_\alpha\cdot\left[\int(N_{I,\beta}D_{\alpha\beta\gamma}N_{J,\eta} + \kappa GhN_IN_J\delta_{\alpha\gamma})d\Omega\right]\cdot\phi_\gamma
\end{equation}
\begin{equation}
\boxed{k_{I\alpha,J\gamma}^{\varphi\varphi} = \int(N_{I,\beta}D_{\alpha\beta\gamma\eta}N_{J,\eta} + \kappa GhN_IN_J\delta_{\alpha\gamma})d\Omega}
\end{equation}
（其中 $\kappa$ 为剪切修正系数，通常薄板或中厚板取 $5/6$）

几何刚度矩阵子块展开如下：展开后的对称矩阵形式为：
\begin{equation}
k_g = \begin{bmatrix}
k_g^{ww} & & \\
& k_g^{\varphi_1\varphi_1} & k_g^{\varphi_1\varphi_2} \\
& k_g^{\varphi_2\varphi_1} & k_g^{\varphi_2\varphi_2}
\end{bmatrix}
\end{equation}

对于 $k_g^{ww}$，
\begin{equation}
k_g^{ww} = \int w_{,\alpha}(h\bar{\sigma}_{\alpha\beta})w_{,\beta}d\Omega
\end{equation}
\begin{equation}
k_g^{ww} = \int(N_{I,\alpha}w_I)\cdot(h\bar{\sigma}_{\alpha\beta})\cdot(N_{J,\beta}w_J)d\Omega
\end{equation}
\begin{equation}
k_g^{ww} = w_I\cdot\left[\int h\bar{\sigma}_{\alpha\beta}N_{I,\alpha}N_{J,\beta}d\Omega\right]\cdot w_J
\end{equation}
\begin{equation}
\boxed{k_{g,IJ}^{ww} = \int h\bar{\sigma}_{\alpha\beta}N_{I,\alpha}N_{J,\beta}d\Omega}
\end{equation}

对于 $k_g^{\varphi_\alpha\varphi_\gamma}$，
\begin{equation}
k_{g,\alpha\gamma}^{\varphi\varphi} = \int \varphi_{,\beta}\left(\frac{h^3}{12}\bar{\sigma}_{\beta}\right)\varphi_{,\gamma}d\Omega
\end{equation}
\begin{equation}
k_{g,I\alpha,J\gamma}^{\varphi\varphi} = \int(N_{J,\beta}\varphi_{\alpha,\alpha})\cdot\left(\frac{h^3}{12}\bar{\sigma}_{\beta}\right)\cdot(N_{I,\varphi_{\gamma,\gamma}})d\Omega
\end{equation}
\begin{equation}
k_{g,I\alpha,J\gamma}^{\varphi\varphi} = \phi_\alpha\cdot\left[\int\frac{h^3}{12}\bar{\sigma}_{\beta}N_{I,\beta}N_{J,\gamma}d\Omega\right]\cdot\phi_\gamma
\end{equation}
\begin{equation}
\boxed{k_{g,I\alpha,J\gamma}^{\varphi\varphi} = \int\frac{h^3}{12}\bar{\sigma}_{\beta}N_{I,\beta}N_{J,\gamma}d\Omega}
\end{equation}

将上述所有子块合并，针对中厚板面外屈曲模式，最终的广义特征值矩阵方程为：
\begin{equation}
\begin{pmatrix}
k_{ww} & k_{w\varphi_1} & k_{w\varphi_2} \\
k_{\varphi_1 w} & k_{\varphi_1\varphi_1} & k_{\varphi_1\varphi_2} \\
k_{\varphi_2 w} & k_{\varphi_2\varphi_1} & k_{\varphi_2\varphi_2}
\end{pmatrix}
+ \lambda
\begin{pmatrix}
k_g^{ww} & & \\
& \frac{h^2}{12}k_g^{ww} & \\
& & \frac{h^2}{12}k_g^{ww}
\end{pmatrix}
\begin{pmatrix} w \\ \varphi_1 \\ \varphi_2 \end{pmatrix}
= \begin{pmatrix} 0 \\ 0 \\ 0 \end{pmatrix}
\end{equation}

% TODO: 待补充——混合离散格式的LBB稳定性分析
% 如需补充，应结合第二章的约束比理论，说明本格式的剪切应力场插值阶数满足inf-sup条件
% 可参考徐洋涛论文第3章的写法：独立成节，含特征值问题推导和约束比数值验证


% 方板屈曲数值算例
% 包含：CCCC, SSSS, CSCS, FSCS, FSSS, SCSC 六种边界条件

\section{数值算例}

核心要点：针对单轴均匀压缩载荷，系统评估混合离散格式在极薄极限下的临界载荷计算能力。以下给出六种边界条件下、不同厚跨比（v01–v12）的屈曲临界荷载系数对比图。

计算模型如图所示，方板求解域为 $\Omega = [-0.5, 0.5] \times [-0.5, 0.5]$，采用均匀节点离散。四条边分别编号为 1（底边）、2（右边）、3（顶边）、4（左边），通过组合不同的边界约束条件（固支 C、简支 S、自由 F）构成六种典型工况。边界箭头表示各边的位移约束方向。

\begin{figure}[H]
    \centering
    \includegraphics[width=0.7\textwidth]{plate_model_generated.png}
    \caption{方板计算模型与边界条件示意图}
    \label{CH4_plate_model}
\end{figure}

\textbf{材料与几何参数}：本文数值算例采用无量纲材料参数，具体取值如表所示。方板边长 $a=b=1.0$，板厚 $h=10^{-2}$，厚跨比 $h/a=1/100$，属于薄板范畴。剪切修正系数取 Mindlin 板理论常用值 $\kappa=5/6$。

\begin{table}[H]
\centering
\caption{方板屈曲算例材料与几何参数}
\label{CH4_params}
\begin{tabular}{lll}
\hline
\textbf{参数} & \textbf{符号} & \textbf{取值} \\
\hline
弹性模量 & $E$ & $1.0$ \\
泊松比 & $\nu$ & $0.3$ \\
密度 & $\rho$ & $1.0$ \\
板厚 & $h$ & $10^{-2}$ \\
板边长 & $a=b$ & $1.0$ \\
剪切修正系数 & $\kappa$ & $5/6$ \\
弯曲刚度 & $D^b$ & $\approx 9.16\times10^{-8}$ \\
剪切刚度 & $D^s$ & $\approx 3.21\times10^{-3}$ \\
预加应力 & $\sigma_{11},\sigma_{22},\sigma_{12}$ & $1.0,0,0$ \\
\hline
\end{tabular}
\end{table}

\textbf{网格离散}：数值算例采用结构化网格离散，四边形单元（Quad4）网格规模为 $25\times25$ 至 $35\times35$，三角形单元（Tri3）网格规模对应约 $392$ 至 $800$ 个单元。典型网格如图所示。

\begin{figure}[H]
    \centering
    \begin{tabular}{cc}
    \includegraphics[width=0.45\textwidth]{meshes/mesh_quad4_st_q_25.png} &
    \includegraphics[width=0.45\textwidth]{meshes/mesh_tri3_st_t_15.png} \\
    (a) Quad4 单元 ($25\times25=625$ 单元) & (b) Tri3 单元 ($392$ 单元)
    \end{tabular}
    \caption{方板结构化网格离散示意图}
    \label{CH4_mesh}
\end{figure}

\begin{table}[H]
\centering
\caption{单轴受压方板屈曲临界荷载系数对比（三角形 Tri3 网格，$35\times35$，$t/b=0.01$）}
\label{CH4_compare_table}
\begin{tabular}{cccccccc}
\hline
边界条件 & 标准解$^{[1]}$ & 文献[Bui] & MF & MF2 & FEM & FEM缩减积分 & FEM三变量 \\
\hline
CCCC & 10.0700 & 10.1080 & 10.0568 & 10.1040 & 8.2847 & 8.2847 & 10.0842 \\
SSSS & 4.0000 & 4.0050 & 3.9921 & 3.9987 & 3.5966 & 3.5966 & 4.0081 \\
CSCS & 7.6900 & 7.7020 & 7.6975 & 7.6617 & 6.2903 & 6.2903 & 7.7069 \\
SCSC & 6.7500 & 6.7800 & 6.7292 & 6.7620 & 5.8397 & 5.8397 & 6.7392 \\
FSCS & 1.7000 & 1.6930 & 1.6497 & 1.6528 & 1.5470 & 1.5470 & 1.6547 \\
FSSS & 1.4400 & 1.4390 & 1.3993 & 1.4019 & 1.3224 & 1.3224 & 1.4032 \\
\hline
\end{tabular}
\end{table}

\begin{figure}[H]
    \centering
    \begin{tabular}{cc}
    \includegraphics[width=0.48\textwidth]{convergence/conv_tri3_CCCC.png} &
    \includegraphics[width=0.48\textwidth]{convergence/conv_tri3_SSSS.png} \\
    (a) 四边固支 CCCC & (b) 四边简支 SSSS
    \end{tabular}
    \caption{三角形（Tri3）网格下各方法屈曲临界荷载系数的对数相对误差收敛曲线}
    \label{CH4_convergence}
\end{figure}


\subsection{方板边界条件屈曲分析}

核心要点：计算四边固支方板在不同厚跨比下的屈曲载荷系数，展示在极薄极限下完全不锁死的数值优越性。

\begin{figure}[H]
    \centering
    模態  1\\[1mm]
    \includegraphics[width=\linewidth]{combined/layout_1x6_v9/CCCC_v01_1x6.png}\\[3mm]
    模態  2\\[1mm]
    \includegraphics[width=\linewidth]{combined/layout_1x6_v9/CCCC_v02_1x6.png}\\[3mm]
    模態  3\\[1mm]
    \includegraphics[width=\linewidth]{combined/layout_1x6_v9/CCCC_v03_1x6.png}\\[3mm]
    模態  4\\[1mm]
    \includegraphics[width=\linewidth]{combined/layout_1x6_v9/CCCC_v04_1x6.png}\\[3mm]
    模態  5\\[1mm]
    \includegraphics[width=\linewidth]{combined/layout_1x6_v9/CCCC_v05_1x6.png}\\[3mm]
    模態  6\\[1mm]
    \includegraphics[width=\linewidth]{combined/layout_1x6_v9/CCCC_v06_1x6.png}\\[3mm]
    \caption{四边固支（CCCC）方板屈曲临界荷载系数对比（v01–v12）（模态1–6）}
    \label{CH4_buckling_CCCC_a}
\end{figure}
\begin{figure}[H]
    \centering
    模態  7\\[1mm]
    \includegraphics[width=\linewidth]{combined/layout_1x6_v9/CCCC_v07_1x6.png}\\[3mm]
    模態  8\\[1mm]
    \includegraphics[width=\linewidth]{combined/layout_1x6_v9/CCCC_v08_1x6.png}\\[3mm]
    模態  9\\[1mm]
    \includegraphics[width=\linewidth]{combined/layout_1x6_v9/CCCC_v09_1x6.png}\\[3mm]
    模態 10\\[1mm]
    \includegraphics[width=\linewidth]{combined/layout_1x6_v9/CCCC_v10_1x6.png}\\[3mm]
    模態 11\\[1mm]
    \includegraphics[width=\linewidth]{combined/layout_1x6_v9/CCCC_v11_1x6.png}\\[3mm]
    模態 12\\[1mm]
    \includegraphics[width=\linewidth]{combined/layout_1x6_v9/CCCC_v12_1x6.png}\\[3mm]
    \caption{四边固支（CCCC）方板屈曲临界荷载系数对比（v01–v12）（模态7–12）}
    \label{CH4_buckling_CCCC_b}
\end{figure}

核心要点：分析四边简支方板的屈曲模态切换特征，探讨最优约束比的收敛规律。

\begin{figure}[H]
    \centering
    模態  1\\[1mm]
    \includegraphics[width=\linewidth]{combined/layout_1x6_v9/SSSS_v01_1x6.png}\\[3mm]
    模態  2\\[1mm]
    \includegraphics[width=\linewidth]{combined/layout_1x6_v9/SSSS_v02_1x6.png}\\[3mm]
    模態  3\\[1mm]
    \includegraphics[width=\linewidth]{combined/layout_1x6_v9/SSSS_v03_1x6.png}\\[3mm]
    模態  4\\[1mm]
    \includegraphics[width=\linewidth]{combined/layout_1x6_v9/SSSS_v04_1x6.png}\\[3mm]
    模態  5\\[1mm]
    \includegraphics[width=\linewidth]{combined/layout_1x6_v9/SSSS_v05_1x6.png}\\[3mm]
    模態  6\\[1mm]
    \includegraphics[width=\linewidth]{combined/layout_1x6_v9/SSSS_v06_1x6.png}\\[3mm]
    \caption{四边简支（SSSS）方板屈曲临界荷载系数对比（v01–v12）（模态1–6）}
    \label{CH4_buckling_SSSS_a}
\end{figure}
\begin{figure}[H]
    \centering
    模態  7\\[1mm]
    \includegraphics[width=\linewidth]{combined/layout_1x6_v9/SSSS_v07_1x6.png}\\[3mm]
    模態  8\\[1mm]
    \includegraphics[width=\linewidth]{combined/layout_1x6_v9/SSSS_v08_1x6.png}\\[3mm]
    模態  9\\[1mm]
    \includegraphics[width=\linewidth]{combined/layout_1x6_v9/SSSS_v09_1x6.png}\\[3mm]
    模態 10\\[1mm]
    \includegraphics[width=\linewidth]{combined/layout_1x6_v9/SSSS_v10_1x6.png}\\[3mm]
    模態 11\\[1mm]
    \includegraphics[width=\linewidth]{combined/layout_1x6_v9/SSSS_v11_1x6.png}\\[3mm]
    模態 12\\[1mm]
    \includegraphics[width=\linewidth]{combined/layout_1x6_v9/SSSS_v12_1x6.png}\\[3mm]
    \caption{四边简支（SSSS）方板屈曲临界荷载系数对比（v01–v12）（模态7–12）}
    \label{CH4_buckling_SSSS_b}
\end{figure}

\begin{figure}[H]
    \centering
    模態  1\\[1mm]
    \includegraphics[width=\linewidth]{combined/layout_1x6_v9/CSCS_v01_1x6.png}\\[3mm]
    模態  2\\[1mm]
    \includegraphics[width=\linewidth]{combined/layout_1x6_v9/CSCS_v02_1x6.png}\\[3mm]
    模態  3\\[1mm]
    \includegraphics[width=\linewidth]{combined/layout_1x6_v9/CSCS_v03_1x6.png}\\[3mm]
    模態  4\\[1mm]
    \includegraphics[width=\linewidth]{combined/layout_1x6_v9/CSCS_v04_1x6.png}\\[3mm]
    模態  5\\[1mm]
    \includegraphics[width=\linewidth]{combined/layout_1x6_v9/CSCS_v05_1x6.png}\\[3mm]
    模態  6\\[1mm]
    \includegraphics[width=\linewidth]{combined/layout_1x6_v9/CSCS_v06_1x6.png}\\[3mm]
    \caption{固支-简支交替（CSCS）方板屈曲临界荷载系数对比（v01–v12）（模态1–6）}
    \label{CH4_buckling_CSCS_a}
\end{figure}
\begin{figure}[H]
    \centering
    模態  7\\[1mm]
    \includegraphics[width=\linewidth]{combined/layout_1x6_v9/CSCS_v07_1x6.png}\\[3mm]
    模態  8\\[1mm]
    \includegraphics[width=\linewidth]{combined/layout_1x6_v9/CSCS_v08_1x6.png}\\[3mm]
    模態  9\\[1mm]
    \includegraphics[width=\linewidth]{combined/layout_1x6_v9/CSCS_v09_1x6.png}\\[3mm]
    模態 10\\[1mm]
    \includegraphics[width=\linewidth]{combined/layout_1x6_v9/CSCS_v10_1x6.png}\\[3mm]
    模態 11\\[1mm]
    \includegraphics[width=\linewidth]{combined/layout_1x6_v9/CSCS_v11_1x6.png}\\[3mm]
    模態 12\\[1mm]
    \includegraphics[width=\linewidth]{combined/layout_1x6_v9/CSCS_v12_1x6.png}\\[3mm]
    \caption{固支-简支交替（CSCS）方板屈曲临界荷载系数对比（v01–v12）（模态7–12）}
    \label{CH4_buckling_CSCS_b}
\end{figure}

\begin{figure}[H]
    \centering
    模態  1\\[1mm]
    \includegraphics[width=\linewidth]{combined/layout_1x6_v9/FSCS_v01_1x6.png}\\[3mm]
    模態  2\\[1mm]
    \includegraphics[width=\linewidth]{combined/layout_1x6_v9/FSCS_v02_1x6.png}\\[3mm]
    模態  3\\[1mm]
    \includegraphics[width=\linewidth]{combined/layout_1x6_v9/FSCS_v03_1x6.png}\\[3mm]
    模態  4\\[1mm]
    \includegraphics[width=\linewidth]{combined/layout_1x6_v9/FSCS_v04_1x6.png}\\[3mm]
    模態  5\\[1mm]
    \includegraphics[width=\linewidth]{combined/layout_1x6_v9/FSCS_v05_1x6.png}\\[3mm]
    模態  6\\[1mm]
    \includegraphics[width=\linewidth]{combined/layout_1x6_v9/FSCS_v06_1x6.png}\\[3mm]
    \caption{自由-简支-固支-简支（FSCS）方板屈曲临界荷载系数对比（v01–v12）（模态1–6）}
    \label{CH4_buckling_FSCS_a}
\end{figure}
\begin{figure}[H]
    \centering
    模態  7\\[1mm]
    \includegraphics[width=\linewidth]{combined/layout_1x6_v9/FSCS_v07_1x6.png}\\[3mm]
    模態  8\\[1mm]
    \includegraphics[width=\linewidth]{combined/layout_1x6_v9/FSCS_v08_1x6.png}\\[3mm]
    模態  9\\[1mm]
    \includegraphics[width=\linewidth]{combined/layout_1x6_v9/FSCS_v09_1x6.png}\\[3mm]
    模態 10\\[1mm]
    \includegraphics[width=\linewidth]{combined/layout_1x6_v9/FSCS_v10_1x6.png}\\[3mm]
    模態 11\\[1mm]
    \includegraphics[width=\linewidth]{combined/layout_1x6_v9/FSCS_v11_1x6.png}\\[3mm]
    模態 12\\[1mm]
    \includegraphics[width=\linewidth]{combined/layout_1x6_v9/FSCS_v12_1x6.png}\\[3mm]
    \caption{自由-简支-固支-简支（FSCS）方板屈曲临界荷载系数对比（v01–v12）（模态7–12）}
    \label{CH4_buckling_FSCS_b}
\end{figure}

\begin{figure}[H]
    \centering
    模態  1\\[1mm]
    \includegraphics[width=\linewidth]{combined/layout_1x6_v9/FSSS_v01_1x6.png}\\[3mm]
    模態  2\\[1mm]
    \includegraphics[width=\linewidth]{combined/layout_1x6_v9/FSSS_v02_1x6.png}\\[3mm]
    模態  3\\[1mm]
    \includegraphics[width=\linewidth]{combined/layout_1x6_v9/FSSS_v03_1x6.png}\\[3mm]
    模態  4\\[1mm]
    \includegraphics[width=\linewidth]{combined/layout_1x6_v9/FSSS_v04_1x6.png}\\[3mm]
    模態  5\\[1mm]
    \includegraphics[width=\linewidth]{combined/layout_1x6_v9/FSSS_v05_1x6.png}\\[3mm]
    模態  6\\[1mm]
    \includegraphics[width=\linewidth]{combined/layout_1x6_v9/FSSS_v06_1x6.png}\\[3mm]
    \caption{自由-简支-简支-简支（FSSS）方板屈曲临界荷载系数对比（v01–v12）（模态1–6）}
    \label{CH4_buckling_FSSS_a}
\end{figure}
\begin{figure}[H]
    \centering
    模態  7\\[1mm]
    \includegraphics[width=\linewidth]{combined/layout_1x6_v9/FSSS_v07_1x6.png}\\[3mm]
    模態  8\\[1mm]
    \includegraphics[width=\linewidth]{combined/layout_1x6_v9/FSSS_v08_1x6.png}\\[3mm]
    模態  9\\[1mm]
    \includegraphics[width=\linewidth]{combined/layout_1x6_v9/FSSS_v09_1x6.png}\\[3mm]
    模態 10\\[1mm]
    \includegraphics[width=\linewidth]{combined/layout_1x6_v9/FSSS_v10_1x6.png}\\[3mm]
    模態 11\\[1mm]
    \includegraphics[width=\linewidth]{combined/layout_1x6_v9/FSSS_v11_1x6.png}\\[3mm]
    模態 12\\[1mm]
    \includegraphics[width=\linewidth]{combined/layout_1x6_v9/FSSS_v12_1x6.png}\\[3mm]
    \caption{自由-简支-简支-简支（FSSS）方板屈曲临界荷载系数对比（v01–v12）（模态7–12）}
    \label{CH4_buckling_FSSS_b}
\end{figure}

\begin{figure}[H]
    \centering
    模態  1\\[1mm]
    \includegraphics[width=\linewidth]{combined/layout_1x6_v9/SCSC_v01_1x6.png}\\[3mm]
    模態  2\\[1mm]
    \includegraphics[width=\linewidth]{combined/layout_1x6_v9/SCSC_v02_1x6.png}\\[3mm]
    模態  3\\[1mm]
    \includegraphics[width=\linewidth]{combined/layout_1x6_v9/SCSC_v03_1x6.png}\\[3mm]
    模態  4\\[1mm]
    \includegraphics[width=\linewidth]{combined/layout_1x6_v9/SCSC_v04_1x6.png}\\[3mm]
    模態  5\\[1mm]
    \includegraphics[width=\linewidth]{combined/layout_1x6_v9/SCSC_v05_1x6.png}\\[3mm]
    模態  6\\[1mm]
    \includegraphics[width=\linewidth]{combined/layout_1x6_v9/SCSC_v06_1x6.png}\\[3mm]
    \caption{简支-固支交替（SCSC）方板屈曲临界荷载系数对比（v01–v12）（模态1–6）}
    \label{CH4_buckling_SCSC_a}
\end{figure}
\begin{figure}[H]
    \centering
    模態  7\\[1mm]
    \includegraphics[width=\linewidth]{combined/layout_1x6_v9/SCSC_v07_1x6.png}\\[3mm]
    模態  8\\[1mm]
    \includegraphics[width=\linewidth]{combined/layout_1x6_v9/SCSC_v08_1x6.png}\\[3mm]
    模態  9\\[1mm]
    \includegraphics[width=\linewidth]{combined/layout_1x6_v9/SCSC_v09_1x6.png}\\[3mm]
    模態 10\\[1mm]
    \includegraphics[width=\linewidth]{combined/layout_1x6_v9/SCSC_v10_1x6.png}\\[3mm]
    模態 11\\[1mm]
    \includegraphics[width=\linewidth]{combined/layout_1x6_v9/SCSC_v11_1x6.png}\\[3mm]
    模態 12\\[1mm]
    \includegraphics[width=\linewidth]{combined/layout_1x6_v9/SCSC_v12_1x6.png}\\[3mm]
    \caption{简支-固支交替（SCSC）方板屈曲临界荷载系数对比（v01–v12）（模态7–12）}
    \label{CH4_buckling_SCSC_b}
\end{figure}


\section{小结}

% TODO: 第四章小结

% NOTE 毕业论文致谢部分
% \include{acknowledgments}
\appendix
\chapter{附录}
附录章节前需添加以下下代码
\begin{lstlisting}
    \appendix
\end{lstlisting}


\chapter{参考格式}
\begin{lstlisting}

\documentclass[internationalmaster]{hquThesis}-hquThesis是模板文件,包含文本格式,字体设置,文献格式.....
\usepackage{amsmath,amsfonts,amssymb,textcomp}-公式所需
\usepackage{booktabs,multirow,tabularx,float}-表格所需
\hypersetup{colorlinks=true,linkcolor=black,citecolor=black}-设置超链接格式
\titleZh{XXX}-中文标题
\titleEn{XXX}-英文标题
\id{}-学号
\authorZh{何祖彦}-作者姓名
\supervisorZh{吴俊超}{副教授}-指导教师+职称
\cosupervisorZh{XXX}{XXX}-合作教师+职称
\practicevisorZh{XXX}{XXX}-实践教师+职称
\departmentZh{土木工程学院}
\fieldZh{结构体系创新与工程应用}
\majortype{工程硕士}
\major{土木水利}
\coverdate{二〇二四年三月XX日}
-封面设置
\begin{document}-正文开始

\makecover
\include{decision}-答辩委员会
\include{declaration}-独创性说明
\frontmatter
\include{abstract}-摘要
\tableofcontents
\mainmatter
\include{Introduction}-正文章节
\include{XXX}
\include{XXXX}
\backmatter-页眉另起设置
\bibliography{references}-引用参考文献条目(.bib)
\include{acknowledgments}-致谢
\appendix-附录部分
\include{appendixA}-附录章节
\include{appendixB}-附录章节
\include{cv}-个人简历

\end{document}-结束正文
\end{lstlisting}


\backmatter


\bibliography{references}

% NOTE 毕业论文个人简历部分
\include{cv}
\include{ai}

\end{document}
